\documentclass[10pt,a4paper]{amsart}
\usepackage{amsmath,amssymb,amsthm,graphicx}
\usepackage{verbatim}
\newenvironment{reference}{\comment}{\endcomment}
%\newenvironment{reference}{}{}
\newenvironment{slogan}{\comment}{\endcomment}
\newenvironment{history}{\comment}{\endcomment}

% For commutative diagrams we use Xy-pic
\usepackage[all]{xy}

% We use 2cell for 2-commutative diagrams.
\xyoption{2cell}
% Theorem environments.
%
\theoremstyle{plain}
\newtheorem{theorem}[subsection]{Theorem}
\newtheorem{proposition}[subsection]{Proposition}
\newtheorem{lemma}[subsection]{Lemma}

\theoremstyle{definition}
\newtheorem{definition}[subsection]{Definition}
\newtheorem{example}[subsection]{Example}
\newtheorem{exercise}[subsection]{Exercise}
\newtheorem{situation}[subsection]{Situation}

\theoremstyle{remark}
\newtheorem{remark}[subsection]{Remark}
\newtheorem{remarks}[subsection]{Remarks}

\numberwithin{equation}{subsection}

% Macros
%
\def\lim{\mathop{\mathrm{lim}}\nolimits}
\def\colim{\mathop{\mathrm{colim}}\nolimits}
\def\Spec{\mathop{\mathrm{Spec}}}
\def\Hom{\mathop{\mathrm{Hom}}\nolimits}
\def\Ext{\mathop{\mathrm{Ext}}\nolimits}
\def\SheafHom{\mathop{\mathcal{H}\!\mathit{om}}\nolimits}
\def\SheafExt{\mathop{\mathcal{E}\!\mathit{xt}}\nolimits}
\def\Sch{\mathit{Sch}}
\def\Mor{\mathop{\mathrm{Mor}}\nolimits}
\def\Ob{\mathop{\mathrm{Ob}}\nolimits}
\def\Sh{\mathop{\mathit{Sh}}\nolimits}
\def\NL{\mathop{N\!L}\nolimits}
\def\CH{\mathop{\mathrm{CH}}\nolimits}
\def\proetale{{pro\text{-}\acute{e}tale}}
\def\etale{{\acute{e}tale}}
\def\QCoh{\mathit{QCoh}}
\def\Ker{\mathop{\mathrm{Ker}}}
\def\Im{\mathop{\mathrm{Im}}}
\def\Coker{\mathop{\mathrm{Coker}}}
\def\Coim{\mathop{\mathrm{Coim}}}

% Boxtimes
%
\DeclareMathSymbol{\boxtimes}{\mathbin}{AMSa}{"02}

%
% Macros for moduli stacks/spaces
%
\def\QCohstack{\mathcal{QC}\!\mathit{oh}}
\def\Cohstack{\mathcal{C}\!\mathit{oh}}
\def\Spacesstack{\mathcal{S}\!\mathit{paces}}
\def\Quotfunctor{\mathrm{Quot}}
\def\Hilbfunctor{\mathrm{Hilb}}
\def\Curvesstack{\mathcal{C}\!\mathit{urves}}
\def\Polarizedstack{\mathcal{P}\!\mathit{olarized}}
\def\Complexesstack{\mathcal{C}\!\mathit{omplexes}}
% \Pic is the operator that assigns to X its picard group, usage \Pic(X)
% \Picardstack_{X/B} denotes the Picard stack of X over B
% \Picardfunctor_{X/B} denotes the Picard functor of X over B
\def\Pic{\mathop{\mathrm{Pic}}\nolimits}
\def\Picardstack{\mathcal{P}\!\mathit{ic}}
\def\Picardfunctor{\mathrm{Pic}}
\def\Deformationcategory{\mathcal{D}\!\mathit{ef}}
\usepackage{cleveref}
\crefname{originallemma}{original lemma}{original lemmas}
\Crefname{originallemma}{Original lemma}{Original lemmas}
\makeatletter
\newlabel{lemma-relocalize-morphism}{{04M4}{source}{}{sourcepointer-1}{}}
\newlabel{lemma-relocalize-morphism@cref}{{[originallemma][0][]04M4}{[1][0][]source}{}{}{}}
\makeatother
\crefname{originalreference}{original reference}{original references}
\Crefname{originalreference}{Original reference}{Original references}
\makeatletter
\newlabel{sites-lemma-relocalize-morphism}{{04IN}{source}{}{sourcepointer-2}{}}
\newlabel{sites-lemma-relocalize-morphism@cref}{{[originalreference][0][]04IN}{[1][0][]source}{}{}{}}
\makeatother
\crefname{originalreference}{original reference}{original references}
\Crefname{originalreference}{Original reference}{Original references}
\makeatletter
\newlabel{sites-lemma-essential-image-j-shriek}{{00Y1}{source}{}{sourcepointer-3}{}}
\newlabel{sites-lemma-essential-image-j-shriek@cref}{{[originalreference][0][]00Y1}{[1][0][]source}{}{}{}}
\makeatother
\crefname{originalreference}{original reference}{original references}
\Crefname{originalreference}{Original reference}{Original references}
\makeatletter
\newlabel{sites-modules-lemma-localize-compare}{{04J3}{source}{}{sourcepointer-4}{}}
\newlabel{sites-modules-lemma-localize-compare@cref}{{[originalreference][0][]04J3}{[1][0][]source}{}{}{}}
\makeatother
\crefname{originalreference}{original reference}{original references}
\Crefname{originalreference}{Original reference}{Original references}
\makeatletter
\newlabel{etale-cohomology-lemma-faithful}{{04LW}{source}{}{sourcepointer-5}{}}
\newlabel{etale-cohomology-lemma-faithful@cref}{{[originalreference][0][]04LW}{[1][0][]source}{}{}{}}
\makeatother
\title{2-isomorphic morphisms of algebraic spaces are equal}
\author{The Stacks Project Authors}
\date{stacks-project:spaces-properties.tex:local-snapshot-20261012; GNU FDL 1.2 or later}
\begin{document}
\maketitle
\noindent\textit{Editorial note.} All delivered mathematical text is the author's own, copied verbatim from the pinned Stacks Project chapter file spaces-properties.tex; no statement or proof step was rewritten, reordered or removed.\par\medskip
\noindent\textit{Editorial note.} Span recomputation by label pairing (LANE\_BRIEF 10.2): in this chapter the author's proof environment follows the statement environment as a standalone block; the spec covers the complete environment pair (statement span, proof span) as recorded in the pool candidate.\par\medskip
\noindent\textit{Editorial note.} Occupation scan: combined claim\_key + statement\_anchor + generic-line scan over all registration files (active and withdrawn) found 13 S10-geom published cards for spaces-properties.tex (prop-occupied.json, 2026-10-14); those 13 labels are excluded from this batch by design. For the 77 free labels: 0 occupied. The registry mentions field (7557 files) and focused-seeds (645) both point to the same S10-geom cards (dual-source check), no other occupation.\par\medskip
\ifdefined\hypertarget\hypertarget{sourcepointer-1}{}\fi
\noindent\textit{Source reference.} Original lemma 04M4: sources/stacks-chow/tags.txt, source lines 6030--6030. Cross-chapter standard prerequisite cited by the author inside the counted span; it is used as a pointer and its own proof is not part of this item.\par\medskip
\ifdefined\hypertarget\hypertarget{sourcepointer-2}{}\fi
\noindent\textit{Source reference.} Original reference 04IN: sources/stacks-chow/tags.txt, source lines 5909--5909. Cross-chapter standard prerequisite cited by the author inside the counted span; it is used as a pointer and its own proof is not part of this item.\par\medskip
\ifdefined\hypertarget\hypertarget{sourcepointer-3}{}\fi
\noindent\textit{Source reference.} Original reference 00Y1: sources/stacks-chow/tags.txt, source lines 1239--1239. Cross-chapter standard prerequisite cited by the author inside the counted span; it is used as a pointer and its own proof is not part of this item.\par\medskip
\ifdefined\hypertarget\hypertarget{sourcepointer-4}{}\fi
\noindent\textit{Source reference.} Original reference 04J3: sources/stacks-chow/tags.txt, source lines 5924--5924. Cross-chapter standard prerequisite cited by the author inside the counted span; it is used as a pointer and its own proof is not part of this item.\par\medskip
\ifdefined\hypertarget\hypertarget{sourcepointer-5}{}\fi
\noindent\textit{Source reference.} Original reference 04LW: sources/stacks-chow/tags.txt, source lines 6022--6022. Cross-chapter standard prerequisite cited by the author inside the counted span; it is used as a pointer and its own proof is not part of this item.\par\medskip
\begin{lemma}
\label{lemma-faithful}
Let $S$ be a scheme.
Let $X$, $Y$ be algebraic spaces over $S$.
Any two morphisms $a, b : X \to Y$ of algebraic spaces over $S$
for which there exists a $2$-isomorphism
$(a_{small}, a^\sharp) \cong (b_{small}, b^\sharp)$
in the $2$-category of ringed topoi are equal.
\end{lemma}

\begin{proof}
Let $t : a_{small}^{-1} \to b_{small}^{-1}$ be the $2$-isomorphism.
We may equivalently think of $t$ as a transformation
$t : a_{spaces, \etale}^{-1} \to b_{spaces, \etale}^{-1}$
since there is not difference between sheaves on $X_\etale$
and sheaves on $X_{spaces, \etale}$.
Choose a commutative diagram
$$
\xymatrix{
U \ar[d]_p \ar[r]_\alpha  & V \ar[d]^q \\
X \ar[r]^a & Y
}
$$
where $U$ and $V$ are schemes, and $p$ and $q$ are surjective \'etale.
Consider the diagram
$$
\xymatrix{
h_U \ar[r]_-\alpha \ar@{=}[d] & a_{spaces, \etale}^{-1}h_V \ar[d]^t \\
h_U \ar@{..>}[r] & b_{spaces, \etale}^{-1}h_V
}
$$
Since the sheaf $b_{spaces, \etale}^{-1}h_V$ is isomorphic to
$h_{V \times_{Y, b} X}$ we see that the dotted arrow comes from a
morphism of schemes
$\beta : U \to V$ fitting into a commutative diagram
$$
\xymatrix{
U \ar[d]_p \ar[r]_\beta  & V \ar[d]^q \\
X \ar[r]^b & Y
}
$$
We claim that there exists a sequence of $2$-isomorphisms
\begin{align*}
(\alpha_{small}, \alpha^\sharp)
& \cong
(\alpha_{spaces, \etale}, \alpha^\sharp) \\
& \cong
(a_{spaces, \etale, c}, a_c^\sharp) \\
& \cong
(b_{spaces, \etale, d}, b_d^\sharp) \\
& \cong
(\beta_{spaces, \etale}, \beta^\sharp) \\
& \cong
(\beta_{small}, \beta^\sharp)
\end{align*}
The first and the last $2$-isomorphisms come from the identifications
between sheaves on $U_{spaces, \etale}$ and sheaves on
$U_\etale$ and similarly for $V$. The second and fourth
$2$-isomorphisms are those of
Lemma \ref{lemma-relocalize-morphism}
with $c : U \to X \times_{a, Y} V$ induced by $\alpha$ and
$d : U \to X \times_{b, Y} V$ induced by $\beta$.
The middle $2$-isomorphism comes from the transformation $t$.
Namely, the functor $a_{spaces, \etale, c}^{-1}$ corresponds
to the functor
$$
(\mathcal{H} \to h_V) \longmapsto
(a_{spaces, \etale}^{-1}\mathcal{H}
\times_{a_{spaces, \etale}^{-1}h_V, \alpha}
h_U \to
h_U)
$$
and similarly for $b_{spaces, \etale, d}^{-1}$, see
Sites, Lemma \ref{sites-lemma-relocalize-morphism}.
This uses the identification of sheaves on $Y_{spaces, \etale}/V$
as arrows $(\mathcal{H} \to h_V)$ in $\Sh(Y_{spaces, \etale})$
and similarly for $U/X$, see
Sites, Lemma \ref{sites-lemma-essential-image-j-shriek}.
Via this identification the structure sheaf $\mathcal{O}_V$ corresponds to the
pair $(\mathcal{O}_Y \times h_V \to h_V)$ and similarly
for $\mathcal{O}_U$, see
Modules on Sites,
Lemma \ref{sites-modules-lemma-localize-compare}.
Since $t$ switches $\alpha$ and $\beta$ we see that $t$ induces an isomorphism
$$
t :
a_{spaces, \etale}^{-1}\mathcal{H}
\times_{a_{spaces, \etale}^{-1}h_V, \alpha}
h_U
\longrightarrow
b_{spaces, \etale}^{-1}\mathcal{H}
\times_{b_{spaces, \etale}^{-1}h_V, \beta}
h_U
$$
over $h_U$ functorially in $(\mathcal{H} \to h_V)$. Also, $t$ is compatible
with $a_c^\sharp$ and $b_d^\sharp$ as $t$ is
compatible with $a^\sharp$ and $b^\sharp$ by our description
of the structure sheaves $\mathcal{O}_U$ and $\mathcal{O}_V$
above. Hence, the morphisms of ringed topoi
$(\alpha_{small}, \alpha^\sharp)$ and $(\beta_{small}, \beta^\sharp)$
are $2$-isomorphic. By
\'Etale Cohomology, Lemma \ref{etale-cohomology-lemma-faithful}
we conclude $\alpha = \beta$! Since $p : U \to X$ is a surjection
of sheaves it follows that $a = b$.
\end{proof}

\end{document}
